\textbf{\textit{Soal. }}Diberikan $\triangle ABC$ dengan $AB = AC$. Lingkaran dalam $\triangle ABC$ menyinggung sisi $AB$ dan $BC$ berturut-turut di $P$ dan $Q$. Garis $CP$ memotong lingkaran dalam $\triangle ABC$ di titik $P$ dan $R$. Jika $PQ = 4$ dan $PR = 5$, panjang $QR$ adalah \dots\\
$\text{a.} \dfrac{1}{2}\sqrt{14}\qquad
\text{b.}2\qquad
\text{c.}\dfrac{3}{5}\sqrt{14}\qquad
\text{d.}\dfrac{4}{5}\sqrt{14}\qquad
\text{e.}3$
\begin{proof}[\textbf{Solusi. }] (Saturday, 11 November 2023)  Misalkan $S$ pada garis $CP$ sehingga $BS \parallel PQ$. Perhatikan bahwa $\triangle CPQ \sim \triangle CSB$. Misalkan panjang $CR = x$ maka $\dfrac{PS}{CP} = \dfrac{QB}{CQ} = 1 \implies PS=CP=5+x$. Selanjutnya, $\dfrac{SB}{PQ}=\dfrac{CB}{CQ}=2 \implies SB = 2PQ=8$. 

Kita punya
$\angle RSB = \angle RPQ$. Lalu, berdasarkan Alternate Segment Theorem pada lingkaran $(RPQ)$ dan garis singgung $AP$ ditemukan $\angle SPB= \angle CPA= \angle PQR$. Oleh karena itu didapat $\triangle PRQ \sim \triangle SBP$.
Dari sini akan diperoleh
\begin{align*}
    \dfrac{PS}{SB} = \dfrac{QP}{RP} \implies \dfrac{5+x}{8} = \dfrac{4}{5} \implies x=\dfrac{7}{5}.
\end{align*}
yang berarti $CR=\dfrac{7}{5}$ dan $CP=\dfrac{32}{5}$.

    \begin{center}
    \begin{asy}
        import olympiad;
        import geometry;
        size(200);
        pair A, B, C, I, P, Q, R, S;
        C = (0,0);
        B = (1,0);
        A = (0.5, 2);
        I = incenter(A,B,C);        
        draw(A--B--C--cycle);
        circle inc = incircle(triangle(A,B,C));
        draw(inc); 
        label("$C$", C, SW);
        label("$B$", B, SE);
        label("$A$", A, N);
        Q = foot(I, B, C);
        P = foot(I, A, B);
        pair PR[] = intersectionpoints(line(P,C), inc);
        R = PR[0];
        draw(C--P);
        draw(Q--R);
        label("$P$", P, SE);
        label("$Q$", Q, SE);
        label("$R$", R, N);
        S = rotate(180, P)*C;
        label("$S$", S, N);
        draw(P--S--B);
        draw(P--Q);
        label("$4$", (P+Q)/2, SE);
        label("$8$", (B+S)/2, SE);
        label("$5$", (P+R)/2, NW);
        label("$x$", (C+R)/2, NW);
        label("$x+5$", (P+S)/2, NW);
        draw(circumcircle(triangle(R,Q,B)), dashed);
    \end{asy}
\end{center}

Sekarang dari Alternate Segment Theorem didapat $\angle CPQ = \angle CQR$ maka $\triangle CQR \sim \triangle CPQ$ sehingga 
\begin{align*}
    \dfrac{CQ}{CR} = \dfrac{CP}{CQ} \implies CQ^2 = CR\cdot CP = \dfrac{7\cdot 32}{25} \implies CQ = \dfrac{4}{5}\sqrt{14}
\end{align*}
yang mengakibatkan 
\begin{align*}
    \dfrac{RQ}{PQ} = \dfrac{CQ}{CP} \implies RQ = PQ\cdot\dfrac{CQ}{CP}=4\cdot \dfrac{\frac{4}{5}\sqrt{14}}{\frac{32}{5}} = \dfrac{1}{2}\sqrt{14}.
\end{align*}
\end{proof}